LLMs predictably improve as model size $N$ and training data size $D$
increase~\citep{hestness2017scalinglaws,hoffmann2022empirical,kaplan2020scaling}.
%
Today, state-of-the-art LLMs are trained at computational scales that
leave no scope for hyperparameter (HP) tuning, although it is widely
accepted that good HPs are critical for effective
training~\citep{yang2022mup,wortsman2023small,bi2024deepseek}.

Both theoretical and empirical efforts have sought to address this.
Theoretically, maximal update parameterization ($\mup$) allows the
optimal learning rate $\etaopt$ and initial weight variance $\varopt$
to remain stable when scaling model
width~\citep{yang2020feature,yang2022mup},
% and depth~\citep{bordelon2023depthwise,yang2023tensor,dey2025dont},
enabling a ``tune small and train large'' strategy.
%
Empirically, DeepSeek~LLM~\citep{bi2024deepseek} adopted ``scaling
laws for HPs,'' where optimal batch size $\bopt$ and optimal learning
rate $\etaopt$ are estimated at small scale, and then extrapolated via
a power law fit in total compute FLOPs, $C$. A similar approach was
used in~\citet{kaplan2020scaling}, forecasting $\bopt$ and $\etaopt$
from loss $L$ and model size $N$.

Relying on a unique predicted $\bopt$ is \emph{inflexible}---it
precludes adjusting $B$ for compute/time trade-offs or hardware
constraints.
%
It is also unclear whether $C$, $L$, $N$, or $D$ (or a combination)
best explains scaling.  \citet{bi2024deepseek} noted, ``for models with
the same [$C$] but different model/data allocations, the optimal
parameter space varies slightly.''
%
Also, no comparable study has been done for weight decay $\lambda$.

This paper introduces a flexible, unified approach to HPs, using both
$\mup$ and scaling laws.
%
We fit power laws to losses derived from hundreds of $\mup$-trained
models, focusing on combinations of $\lambda$, $B$, $N$, and $D$.
%
We study both \emph{compute-optimal} and \emph{overtrained} models.
The fewest FLOPs to achieve a loss typically occurs when training at
$\approx$20 tokens-per-parameter (TPP =
$D/N$)~\citep{hoffmann2022empirical,besiroglu2024chinchilla}, but
overtrained models ($>$20~TPP) offer more-efficient
inference~\citep{touvron2023llama}.  We study TPPs from 20 to 1280.

To capture $\lambda$'s interaction with other scaling HPs ($\eta$,
$B$), we model scaling of the AdamW timescale, $\tepoch =
B/(\eta\lambda D)$.
%
\citet{wang2024how} found optimal $\tepoch$ stable with
varying $D$, but we show it obeys a power law in TPP (\cref{fig:hook},
$\leftfig$).
%
This law thus enables accurate estimation of $\lambdaopt$ for any $N$,
$D$, $B$.

Leveraging these better HPs as $B$ scales, we also study optimal batch
size: the $\bopt$ that minimizes loss at a given $N$ and $D$.
%
While $\bopt$ scales as a power law in $C$ \emph{when TPP is fixed}
(\cref{fig:related}, $\leftfig$), our results show this arises from a
more fundamental power-law dependence on $D$ (\cref{fig:hook},
$\middlefig$).

Importantly, increasing $B > \bopt$ can still reduce training time
(fewer steps) and improve hardware utilization.
%
This raises the question: how much \emph{extra} data is needed when
using large $B$?
%
Prior work defines the \emph{critical batch size} $\cbs$ as the point
where training to a target loss requires $2 \times \dmin$, with
rapidly-diminishing returns in training speed
thereafter~\citep{mccandlish2018empirical}.
%
We show $\cbs$ also scales with $D$ (\cref{fig:hook}, $\rightfig$),
not $L$ as suggested in~\citet{kaplan2020scaling} (\cref{fig:related},
$\middlefig$), consistent with recent results
from~\citet{zhang2024how}.

Finally, amid intense competition to advance LLM performance, a key
question is: which $N$, $D$, and $B$ yield the best trade-off between
training speed and compute cost?
%
Using our fit $\cbs$ law, we derive Pareto-optimal solutions to these
competing objectives, and show that small, overtrained models can be
best---offering both faster steps and greater parallelism via larger
$D$ (and thus higher $\cbs$).

Key findings and takeaways are highlighted in the paper. Our main
contributions are:
\begin{itemize}[topsep=1pt,itemsep=1pt]
\item The first large-scale empirical study varying weight decay $\lambda$ across $N$, $D$, and $B$ in LLMs.
\item Showing the AdamW timescale obeys a power law, enabling $\lambdaopt$ for any $N$, $D$, $B$ (\cref{sec:lambda}).
\item A new method for estimating $\cbs$, suitable for any LR schedule or optimizer (\cref{sec:cbs_methods}).
\item Confirmation that both $\bopt$ and $\cbs$ scale as power laws in $D$ (\cref{sec:bcrit}).
\item New methods for selecting $N$, $D$, and $B$ to trade-off training time vs.\@ compute (\cref{sec:balancing}).
\end{itemize}
